\documentclass{situationist-preprint}

\title{A Free-Will Defence Without Libertarian Assumptions}
\msid{REL-02-b}
\msfamily{REL-02}
\msdiscipline{Philosophy of Religion}
\mstype{Companion manuscript}
\mskeywords{free-will defence; compatibilism; libertarianism; problem of evil; independent goodness; Mackie}

\begin{document}
\maketitle

\begin{abstract}
Plantinga's free-will defence is usually thought to require libertarian freedom, since on compatibilism God could have created free creatures determined always to do right. REL-02-a answers the logical problem of evil with a good that is not freedom: creatures' goodness that is not brought about by God. We show that this defence needs neither of the theses that make up libertarianism, that freedom is incompatible with determinism and that we are free in a sense that requires indeterminism. It needs only that some state of the world that God brings about does not by itself settle how creatures act. That condition is consistent with compatibilism, with determinism of creaturely deliberation, and with the account of self-formation on which the agent's history, not an undetermined choice, makes the agent the source.
\end{abstract}

\section{The difficulty}

The free-will defence holds that a world containing free creatures who sometimes do wrong may be better than any world God could create without wrongdoing \citep{Plantinga1974}. \citet{Mackie1955} and \citet{Flew1955} replied that, if freedom is compatible with determinism, God could have created creatures who are free and determined always to choose the good. The reply has force against the defence as Plantinga stated it, since the good the defence appeals to, free creatures acting well, is then available without evil. The usual response is to adopt libertarianism, the conjunction of incompatibilism with the thesis that we are free \citep{Kane1996}. Two results on agency rule that response out as a premise of the defence: ACT-01 holds that free will is compatible with determinism and with indeterminism alike, and ACT-08 holds that we need not be the ultimate source of our actions in the sense of Strawson's Basic Argument \citep{Strawson1994}.

REL-02-a replaces the good to which the defence appeals. The good is not freedom but the goodness of creatures that is not brought about by God (Assumption~4.1 of REL-02-a), and the Diagonal Lemma (Lemma~3.2) shows that no agent brings about the independent version of any state. The present note shows that the replacement removes every libertarian assumption from the defence.

\section{The argument}

Say that God \emph{settles} the conduct of a creature at a world if God brings about, in the sense of Definition~2.1 of REL-02-a, that the creature acts as it does: an exercise of God's power by itself settles that the conduct occurs.

\begin{assumption}[Release]\label{a:release}
There is a state of affairs $r$ that God can bring about, in which creatures exist and act, and which does not by itself settle how they act.
\end{assumption}

\begin{proposition}\label{p:noLib}
Assumption~\ref{a:release} entails neither incompatibilism nor the thesis that creatures are free in a sense that requires indeterminism.
\end{proposition}

\begin{proof}
For incompatibilism: Assumption~\ref{a:release} mentions no notion of freedom, and is satisfied in models in which freedom is analysed as the exercise of a reasons-responsive mechanism that is the agent's own \citep{FischerRavizza1998}, a property that deterministic and indeterministic processes can equally have. For the second thesis: Assumption~\ref{a:release} requires only that God's act not settle the conduct by itself. This holds in a world whose laws are deterministic if God does not bring about its initial conditions in full, and in a world in which deliberation is deterministic if some chance process upstream of deliberation, such as which circumstances arise, is not settled by God. Neither case places indeterminism in the agent's choice.
\end{proof}

\begin{proposition}\label{p:defence}
Given Assumption~\ref{a:release} and Assumption~4.1 of REL-02-a, the consistency result of REL-02-a (Theorem~4.2) holds in a model in which free will is compatibilist and creaturely deliberation is deterministic.
\end{proposition}

\begin{proof}
The model of REL-02-a requires an option, release, whose outcomes include creatures' acting well without God's bringing it about, and whose value exceeds that of every option on which God settles creaturely conduct. Interpret release as $r$. By Proposition~\ref{p:noLib}, $r$ is available in a model of the stated kind. The Diagonal Lemma applies to any agent and state, and its proof uses no premise about freedom. The remainder of the proof of Theorem~4.2 is unchanged.
\end{proof}

Mackie's reply is accepted and disarmed. There are worlds in which God determines creatures who are free in the compatibilist sense to act well, and REL-02-a grants it (Section~5). Those worlds lack the independent good, because in them God settles the creatures' conduct. The defence asks whether the independent good is present, which depends on the relation between creaturely conduct and God's agency, not on what freedom consists in.

\section{Objections}

\begin{objection}
If the creature's conduct is not settled by God but by chance, its goodness is luck, not the creature's own.
\end{objection}

\begin{reply}
Chance enters at the level of what God settles, not at the level of the creature's agency. Given the circumstances that arise, the creature acts from its own reasons-responsive mechanism, and its conduct is its own by the standards of ACT-01. ACT-08 adds that a creature can be the source of its character through a history of self-formation that does not begin from nothing. The independent good is the goodness of creatures who are sources in that sense of a history God did not settle, which is the good to which REL-02-a appeals.
\end{reply}

\begin{objection}
Under determinism, God, by creating the world, brings about everything in it; the independent good is then impossible.
\end{objection}

\begin{reply}
Creation brings about whatever God's creative act by itself settles. If God creates a world whose initial state is fixed by God and whose laws are deterministic, God settles everything in it and the independent good is absent, as the objection says. A defence needs only that some option of the other kind be possible, and it is: God may bring about a state whose further determination is left to processes God does not settle. The defence is consistent with determinism being true in some worlds God could create; it requires only that not every world God could create be of that kind.
\end{reply}

\section{Conclusion}

The good on which REL-02-a rests is the goodness of creatures that God did not bring about. Its availability depends on God's refraining from settling outcomes, not on any thesis about the metaphysics of freedom. The free-will defence, so understood, is available to compatibilists, and it is consistent with the account of free will and sourcehood that ACT-01 and ACT-08 defend.

\begin{thebibliography}{99}\small
\bibitem[Fischer and Ravizza(1998)]{FischerRavizza1998} Fischer, J.~M. and Ravizza, M. (1998). \emph{Responsibility and Control: A Theory of Moral Responsibility}. Cambridge University Press.
\bibitem[Flew(1955)]{Flew1955} Flew, A. (1955). Divine omnipotence and human freedom. In A. Flew and A. MacIntyre (Eds.), \emph{New Essays in Philosophical Theology}. SCM Press.
\bibitem[Kane(1996)]{Kane1996} Kane, R. (1996). \emph{The Significance of Free Will}. Oxford University Press.
\bibitem[Mackie(1955)]{Mackie1955} Mackie, J.~L. (1955). Evil and omnipotence. \emph{Mind}, 64(254), 200--212.
\bibitem[Plantinga(1974)]{Plantinga1974} Plantinga, A. (1974). \emph{God, Freedom, and Evil}. Harper \& Row.
\bibitem[Strawson(1994)]{Strawson1994} Strawson, G. (1994). The impossibility of moral responsibility. \emph{Philosophical Studies}, 75(1--2), 5--24.
\end{thebibliography}

\end{document}
